\subsubsection*{Перестройка несбалансированного конвейера для максимальной скорости. Формулы $K_y(n, r)$, $e(n, r)$}

Исходные данные: несбалансированный конвейер с временами стадий $t_1, \ldots, t_n$, суммарное время работы одной задачи $W = \sum_{i=1}^{n} t_i$, $t_{\max} = \max_i t_i$; задан минимальный квант времени $t_0$.

\paragraph{Почему максимальная скорость достигается при балансировке.}
Время работы конвейера на $r$ задачах:
\[
T_n = \sum_{i=1}^{n} t_i + (r - 1) t_{\max} = W + (r - 1) t_{\max}.
\]
При фиксированных $W$ и $n$ величина $T_n$ минимальна, когда минимален $t_{\max}$. При заданной сумме $W$ минимум максимума достигается при равном распределении: $t_1 = \ldots = t_n = W/n$, т.е. $t_{\max} = W/n$. Значит, \textbf{максимально быстрый конвейер получается выравниванием времени всех стадий} (сбалансированный конвейер).

\paragraph{Учёт минимального кванта $t_0$.}
Время каждой стадии должно быть не меньше $t_0$ и целесообразно брать кратным $t_0$. После балансировки каждой стадии назначается одно и то же время $t$, причём $t \geq t_0$ и $t = k t_0$ для некоторого натурального $k$. Обычно берут $t = \max\bigl(t_0,\, \lceil W/(n t_0) \rceil \cdot t_0\bigr)$ или, если допустимо округление суммы, $t$ так, чтобы $n t \geq W$ и $t \geq t_0$. Для вывода формул достаточно положить $t_i = t$ для всех $i$, где $t$ — общее время одной стадии после перестройки.

\paragraph{Перестройка конвейера.}
\begin{enumerate}
    \item Выровнять загрузку стадий: присвоить всем стадиям одинаковое время $t$ (например, $t = W/n$ при допустимости дробного кванта или $t = \lceil W/(n t_0) \rceil t_0$ при дискретном $t_0$), чтобы $t_{\max} = t$ и суммарная работа сохранялась или не увеличивалась.
    \item При необходимости округлить $t$ до значения, кратного $t_0$ и не меньшего $t_0$.
\end{enumerate}

После такой перестройки конвейер сбалансирован: $t_1 = \ldots = t_n = t$, $W = n t$, $T_1 = r n t$, $T_n = n t + (r - 1) t = (n + r - 1) t$.

\paragraph{Формулы для максимально быстрого (сбалансированного) конвейера.}
\[
\boxed{
T_1(n, r) = r n t, \quad T_n(n, r) = (n + r - 1) t,
}
\]
\[
\boxed{
K_y(n, r) = \frac{T_1}{T_n} = \frac{r n t}{(n + r - 1) t} = \frac{r n}{n + r - 1},
}
\]
\[
\boxed{
e(n, r) = \frac{K_y}{n} = \frac{r}{n + r - 1}.
}
\]

Здесь $t$ — время одной стадии после балансировки ($t \geq t_0$); в формулах для $K_y$ и $e$ оно сокращается, поэтому они зависят только от $n$ и $r$.

\paragraph{Подстановка исходных данных (как в предыдущем вопросе).}
При $n = 4$ и выбранном $t$ (например, $t = t_0$ или $t = W/4$):
\[
K_y(4, r) = \frac{4r}{3 + r}, \qquad e(4, r) = \frac{r}{3 + r}.
\]

Таким образом, перестройка сводится к балансировке конвейера (выравниванию $t_i$ с учётом $t_0$); для полученного конвейера справедливы приведённые формулы $K_y(n, r)$ и $e(n, r)$.
